% ============================================================
\section{Experimental Setup}
\label{sec:experiments}
% ============================================================

We evaluate GNNocRoute-DRL through cycle-accurate simulations using BookSim2~\cite{booksim}, a widely-used NoC simulator. This section describes the experimental methodology, baselines, traffic patterns, and evaluation metrics.

\subsection{Simulation Environment}

Simulation parameters are summarized in Table~\ref{tab:simparams}.

\begin{table}[h]
\centering
\caption{BookSim2 simulation parameters}
\label{tab:simparams}
\begin{tabular}{ll}
\toprule
\textbf{Parameter} & \textbf{Value} \\
\midrule
Simulator & BookSim2 2.0 \\
Topologies & Mesh $4\times4$, Mesh $8\times8$, Torus $4\times4$ \\
Flit size & 32 bytes \\
Buffer depth & 4 flits/VC \\
Virtual channels & 4 \\
Routing computation & 1 cycle \\
Warmup periods & 3 \\
Sample period & 1,000 cycles \\
Simulation length & 100,000 cycles \\
Packet size & 1 flit \\
\midrule
Seeds per config & 5 \\
Total configs & 252 \\
\bottomrule
\end{tabular}
\end{table}

\subsection{Baselines}

We compare GNNocRoute-DRL against four routing algorithms:

\begin{itemize}
    \item \textbf{XY (DOR):} Dimension-order routing, the most widely used deterministic algorithm for mesh NoCs. Packets are routed first along the X dimension, then Y.
    \item \textbf{Adaptive XY/YX:} An adaptive variant that selects between XY and YX paths based on local congestion at the source router.
    \item \textbf{Minimal Adaptive:} Routes packets along any minimal path, with port selection based on buffer occupancy.
    \item \textbf{Valiant:} Randomized routing that sends packets to a random intermediate node before forwarding to the destination, providing load balancing at the cost of non-minimal paths.
\end{itemize}

\subsection{Traffic Patterns}

We evaluate under three synthetic traffic patterns:

\begin{itemize}
    \item \textbf{Uniform Random:} Each node sends packets to a uniformly random destination. Represents balanced workloads.
    \item \textbf{Transpose:} Node $(i,j)$ sends to node $(j,i)$. Creates diagonal traffic with moderate stress on specific links.
    \item \textbf{Hotspot:} $10\%$ of traffic targets a single hotspot node $(3,3)$ in the center. Tests congestion handling under skewed loads.
\end{itemize}

\subsection{Evaluation Metrics}

We measure four metrics:

\begin{itemize}
    \item \textbf{Average packet latency:} Mean number of cycles from packet injection to tail flit reception.
    \item \textbf{Throughput:} Accepted flit rate measured as flits/cycle/node.
    \item \textbf{Energy consumption:} Total energy including router dynamic power, link power, and buffer power.
    \item \textbf{Saturation throughput:} Injection rate at which latency exceeds $3\times$ the zero-load latency.
\end{itemize}

\subsection{DRL Training Setup}

The DRL agent is trained offline using the NoC routing environment (Section~\ref{sec:method}). Training consists of:
\begin{enumerate}
    \item GNN encoder pre-training: 500 episodes on Mesh $4\times4$
    \item DQN training: 500 episodes on Mesh $4\times4$, 300 on Mesh $8\times8$
    \item Transfer learning: Pre-trained weights from $4\times4$ fine-tuned on $8\times8$ for 100 episodes
\end{enumerate}

All training runs on CPU (Intel Xeon, 2.3 GHz) using PyTorch 2.8 and PyTorch Geometric 2.6.
